% One job, cycle by cycle. Edge 1 samples start; done is visible after edge 10.
\begin{tikzpicture}[font=\scriptsize, x=0.62cm, y=0.5cm, line width=0.5pt]
  \def\lo{0} \def\hi{0.7}
  % cycle grid: cycle c spans [c, c+1], edge e sits at x = e
  \foreach \e in {1,...,11} {
    \draw[black!25] (\e,0.9) -- (\e,-12.3);
    \node[black!60] at (\e,1.3) {\e};
  }
  \node[anchor=east, black!60] at (0.6,1.3) {edge};

  % helper: a signal high from x1 to x2 on row y
  \newcommand{\sig}[4]{% name, y, x1, x2
    \node[anchor=east] at (0,#2+0.35) {#1};
    \draw (0,#2) -- (#3,#2) -- (#3,#2+\hi) -- (#4,#2+\hi) -- (#4,#2) -- (12,#2);}

  % clock
  \node[anchor=east] at (0,0.35) {clk};
  \foreach \c in {0,...,11} {
    \draw (\c,0) -- (\c,\hi) -- (\c+0.5,\hi) -- (\c+0.5,0) -- (\c+1,0);
  }
  \sig{start}{-1.5}{0.5}{1.5}
  \sig{load, clear}{-3}{1}{2}
  \sig{run}{-4.5}{2}{9}
  % t bus
  \node[anchor=east] at (0,-5.65) {$t$};
  \draw (0,-6) -- (2,-6) (0,-5.3) -- (2,-5.3);
  \foreach \c in {2,...,8} {
    \pgfmathtruncatemacro{\tv}{\c-2}
    \draw (\c,-6) -- (\c+0.1,-5.3) -- (\c+0.9,-5.3) -- (\c+1,-6) -- (\c+0.1,-6) -- (\c,-5.3);
    \node at (\c+0.5,-5.65) {\tv};
  }
  \draw (9,-6) -- (12,-6) (9,-5.3) -- (12,-5.3);
  \sig{capture}{-7.5}{9}{10}
  \sig{done}{-9}{10}{11}
  % c_flat bus
  \node[anchor=east] at (0,-10.15) {c\_flat};
  \draw (0,-10.5) -- (10,-10.5) (0,-9.8) -- (10,-9.8);
  \draw (10,-10.5) -- (10.1,-9.8) -- (12,-9.8) (10,-9.8) -- (10.1,-10.5) -- (12,-10.5);
  \node at (5,-10.15) {previous result};
  \node at (11.05,-10.15) {C};

  \node[anchor=west, align=left] at (0,-11.9)
    {Compute window: 7 cycles ($3N-2$). Start sampled to done visible: 10 edges ($3N+1$).};
\end{tikzpicture}
